\documentclass[prd,aps,10pt,nofootinbib,twocolumn,superscriptaddress,preprintnumbers,balancelastpage,longbibliography,floatfix]{revtex4-2}

\usepackage{amsmath,amssymb}	
\usepackage{mathtools}
\usepackage{fontawesome}
\usepackage[dvipsnames]{xcolor}
\usepackage{hyperref}
\usepackage{xspace}
\usepackage{fancyhdr}
\usepackage{braket}
\usepackage{graphicx}
\usepackage{blindtext}
\usepackage{nicefrac}
\usepackage{lipsum}
\usepackage{bbold}
\usepackage{tabularx}
\usepackage{multirow}
\usepackage{dcolumn}
\usepackage{makecell}
\usepackage{floatrow}
\usepackage{afterpage}
\usepackage{longtable}
\setlength{\LTcapwidth}{\textwidth}

\input{definitions}

\hypersetup{colorlinks=true,
linkcolor=linkcolor,
citecolor=linkcolor,
urlcolor=linkcolor,
,linktocpage=true
,pdfproducer=medialab}

% define "struts", as suggested by Claudio Beccari in
%    a piece in TeX and TUG News, Vol. 2, 1993.
\newcommand\Tstrut{\rule{0pt}{2.6ex}}         % = `top' strut
\newcommand\Bstrut{\rule[-1.6ex]{0pt}{0pt}}   % = `bottom' strut

\renewcommand{\arraystretch}{1.8}
\newcolumntype{C}[1]{>{\centering\let\newline\\\arraybackslash\hspace{0pt}}m{#1}}

\begin{document}

\preprint{\hfill MIT-CTP/XXXX}

\title{On the morphology of $\gamma$-ray point source populations in the Galactic Center}

\author{Siddharth Mishra-Sharma}
\email{smsharma@mit.edu}
\thanks{ORCID: \href{https://orcid.org/0000-0001-9088-7845}{0000-0001-9088-7845}}
\affiliation{The NSF AI Institute for Artificial Intelligence and Fundamental Interactions}
\affiliation{Center for Theoretical Physics, Massachusetts Institute of Technology, Cambridge, MA 02139, USA}
\affiliation{Department of Physics, Harvard University, Cambridge, MA 02138, USA}
\date{\today}

\begin{abstract}
\end{abstract}

\maketitle

\tableofcontents

\section{Introduction}
\label{sec:intro}

\section{Methodology}
\label{sec:methodology}

\subsection{Datasets and spatial templates}

\subsection{Likelihood and forward model}

\subsection{Inference procedure}

\section{Results}
\label{sec:results}

\section{Discussion and conclusions}
\label{sec:conclusion}

\vspace{.2cm}
%%%%%%%%%%%%%%%

\begin{acknowledgments}
This work is supported by the National Science Foundation under Cooperative Agreement PHY-2019786 (The NSF AI Institute for Artificial Intelligence and Fundamental Interactions, \url{http://iaifi.org/}).
This material is based upon work supported by the U.S. Department of Energy, Office of Science, Office of High Energy Physics of U.S. Department of Energy under grant Contract Number DE-SC0012567.
We thank the \Fermi-LAT Collaboration for making publicly available the $\gamma$-ray data used in this work.
This research has made use of NASA's Astrophysics Data System. 
This research made use of the \texttt{astropy}~\cite{Price-Whelan:2018hus,Robitaille:2013mpa}, \texttt{dynesty}~\cite{Speagle_2020}, \texttt{getdist}~\cite{Lewis:2019xzd}, \texttt{IPython}~\cite{PER-GRA:2007}, \texttt{Jupyter}~\cite{Kluyver2016JupyterN}, \texttt{matplotlib}~\cite{Hunter:2007}, \texttt{MLflow}~\cite{10.1145/3399579.3399867}, \texttt{nflows}~\cite{nflows}, \texttt{NPTFit}~\cite{Mishra-Sharma:2016gis}, \texttt{NPTFit-Sim}~\cite{NPTFit-Sim}, \texttt{NumPy}~\cite{harris2020array}, \texttt{pandas}~\cite{pandas:2010}, \texttt{PyGSP}~\cite{michael_defferrard_2017_1003158}, \texttt{Pyro}~\cite{bingham2019pyro}, \texttt{PyTorch}~\cite{NEURIPS2019_9015}, \texttt{PyTorch Geometric}~\cite{Fey/Lenssen/2019}, \texttt{PyTorch Lightning}~\cite{william_falcon_2020_3828935}, \texttt{seaborn}~\cite{seaborn}, \texttt{sbi}~\cite{tejero-cantero2020sbi}, \texttt{scikit-learn}~\cite{JMLR:v12:pedregosa11a}, \texttt{SciPy}~\cite{2020SciPy-NMeth}, and \texttt{tqdm}~\cite{casper_da_costa_luis_2021_5517697} software packages.
\end{acknowledgments}

\vspace{.5cm}
%%%%%%%%%%%%%%%

\appendix
\section*{Appendix}

\clearpage

\bibliographystyle{apsrev4-1}
\bibliography{fermi-prob}

\end{document}
